% ============================================================================
%  Starting point for a technical talk. Copy this file, keep the shape.
%
%  Build:  ./checkslides.sh slides.tex
%          It refuses to pass while any frame overflows, and prints a style
%          audit. Do not ship a deck it has not passed.
%
%  Speaker copy with notes:
%          change \setbeameroption in slidekit.sty to `show notes`, rebuild.
% ============================================================================
\documentclass[8pt,aspectratio=169]{beamer}
\usepackage{slidekit}

\title{Short noun phrase}
\subtitle{One line saying what this is}
\author{Paper authors \quad\gray{venue, year}}
\date{\gray{Presenter: name \quad occasion, date}}

\begin{document}

\titleframe

% The outline is the ARGUMENT, not a list of nouns. \tableofcontents gives you
% four words in a column down the left of an otherwise empty slide.
%
% But keep it MINIMAL: one row per section, three short columns, nothing else.
% A richer version with a gloss column and small-caps markers was built and
% rejected as cluttered -- it is slide two, the audience has not warmed up, and
% every extra element is one they have to skip. The middle column carries the
% section name; the right column carries the single symbol or number that
% section delivers, so the notation is met once before it matters.
\begin{frame}{Outline}
  \begin{center}
  \renewcommand{\arraystretch}{2.4}
  \begin{tabular}{@{}r@{\hspace{1.4em}}l@{\hspace{2.4em}}l@{}}
    \gray{1} & Problem   & the quantity nobody can define \\
    \gray{2} & Theory    & $\text{the computable bound}$ \\
    \gray{3} & Results   & $\text{the number}$ \\
    \gray{4} & Appraisal & the caveat \\
  \end{tabular}
  \end{center}
\end{frame}

% ===========================================================================
\section{Short section label}
% ===========================================================================

% ---------------------------------------------------------------------------
% THE ONLY FRAME SHAPE. Every frame in the deck looks like this.
%
%   title       a LABEL, 2-4 words. Not a sentence, not a compressed abstract.
%   \takeaway   the finding, ONE line. Written FIRST, before the evidence.
%               If you cannot state it in one line, the frame has no point yet
%               and needs splitting.
%   body        ONE piece of evidence: a figure, or a table, or a derivation.
%               Never a figure and a table showing the same numbers.
%   \note       everything else. Setup, sample sizes, caveats, the sentences
%               you will actually speak. This is where prose lives.
% ---------------------------------------------------------------------------
\begin{frame}{Label}
  \takeaway{The one thing to remember, stated as a claim, on one line.}

  \fig{some_figure}

  \showtakeaway
  \note{Model, dataset, n, protocol. What the axes mean. The caveat you would
    mention if asked. Anything you would have been tempted to put in a
    "how to read this" box on the slide itself.}
\end{frame}

% Optional grey provenance line, measured together with the box so it can
% never tip the frame over the edge:
%   \takeaway[n=5000, single seed]{The claim.}

% Two figures that must be COMPARED go side by side, never stacked. A stacked
% pair wastes the width of a 16:9 frame and defeats comparison. If the paper
% stacked them, split the PDF with croppanel.sh:
%   ./croppanel.sh figures/combined.pdf                     # print the bbox
%   ./croppanel.sh figures/combined.pdf out.pdf x0 y0 x1 y1 # crop
%   -> then LOOK at the .preview.png it writes. Every crop this skill has ever
%      got wrong was one nobody looked at.
\begin{frame}{Label}
  \takeaway{What the comparison shows.}

  \figpair{left_panel}{right_panel}

  \showtakeaway
  \note{}
\end{frame}

% Tables: tint the header row with \headrow, immediately after \toprule.
% Rules are already in the accent colour; black hairlines chop a table into
% bands and compete with the data.
%
%   \begin{tabular}{@{}lc@{}}
%     \toprule
%     \headrow
%     Metric & value \\
%     \midrule
%     ...
%     \bottomrule
%   \end{tabular}

% Two columns: text one side, evidence the other. Keep the text side to three
% short items. If it needs four, the frame is two frames.
%
% [c], not [T], whenever a text column is paired with a FIGURE column. The two
% will never be the same height, and [T] pins the short one to the top where it
% hangs stranded above a hole. Keep [T] only for text against text.
\begin{frame}{Label}
  \takeaway{The claim.}

  \begin{columns}[c]
    \column{0.44\textwidth}
      \begin{itemize}\setlength\itemsep{0.7em}
        \item short phrase
        \item short phrase
        \item short phrase, with \F{the failing thing} marked
      \end{itemize}

    \column{0.52\textwidth}
      \fig{some_figure}
  \end{columns}

  \showtakeaway
  \note{}
\end{frame}

\end{document}
